% ============================== АБСТРАКТ ==============================
\chapter*{Абстракт}
\addcontentsline{toc}{chapter}{Абстракт}
\thispagestyle{plain}

Центровете за данни и облачните инфраструктури концентрират услугите на хиляди организации. Затова те са основна цел на атаките Denial of Service (DoS), на опитите за неоторизиран достъп с Brute Force и на мрежите Botnet. Настоящото изследване изгражда, обучава и експериментално сравнява девет модела за откриване на прониквания (IDS) в мрежовия трафик на център за данни. Използвани са три алгоритъма от различни семейства: \emph{Random Forest}, \emph{метод на опорните вектори с RBF ядро} (SVM) и \emph{градиентно усилване на решаващи дървета} (XGBoost). Те са приложени към три атаки с различен характер: Distributed Denial of Service (DDoS), Brute Force срещу FTP и SSH (BruteForce) и Botnet (Bot).

Изследването използва изключително реални записи от публичния набор CSE-CIC-IDS2018. Наборът е уловен в инфраструктура на Amazon Web Services с 420 машини и 30 сървъра. Изтеглени са \val{rows-read} записа на мрежови потоци с \val{n-features} признака. След почистване (премахване на \val{nan-rows} невалидни записа и \val{dup-rows} дубликата) остават \val{rows-clean}. Векторът на признаците се редуцира в два етапа: филтър по дисперсия (\val{n-features}\,$\to$\,\val{n-filtered}) и подбор на \val{n-selected}-те най-информативни признака по важност в помощен модел, изчислена само върху обучаващите данни. Анализът на главните компоненти (PCA) е контролен вариант. Хиперпараметрите се избират с 5-кратна кръстосана валидация и правилото „една стандартна грешка“, което предпочита по-простия модел и ограничава преобучението. Моделите се оценяват в отделна оперативна фаза върху 30\% невиждани данни по TP, TN, FP, FN, Accuracy, Precision, Recall (Sensitivity), F1 и Specificity.

Всички девет модела постигат Accuracy над 99,5\%. За DDoS тя е \val{RF-DDoS-Accuracy} (RF), \val{SVM-DDoS-Accuracy} (SVM) и \val{XGB-DDoS-Accuracy} (XGBoost). За BruteForce е \val{RF-BruteForce-Accuracy}, \val{SVM-BruteForce-Accuracy} и \val{XGB-BruteForce-Accuracy}, а за Bot е \val{RF-Bot-Accuracy}, \val{SVM-Bot-Accuracy} и \val{XGB-Bot-Accuracy}. Разликата между точността при обучение и при тест е най-много \val{gap-max} процентни пункта, а оценката от кръстосаната валидация съвпада с тестовата, т.е. моделите не са преобучени. Ансамблите от дървета превъзхождат SVM статистически значимо при Bot (тест на McNemar). При естественото съотношение на класовете (\val{nat-share-Bot}--\val{nat-share-DDoS} атаки) Precision остава над 99\%. Контролните експерименти показват, че резултатите не зависят от времевите признаци и от метода на редукция. Допълнителни проверки за приложимост показват, че точността почти не зависи от признаците-отпечатъци на атакуващия инструмент и се запазва при хронологично разделяне за DDoS и BruteForce. При Bot ансамблите пропускат около 10\% от по-късните атаки, а комбинираният детектор XGBoost или SVM ги открива почти всички. Предварителното изследване върху IoT набора CICIoT2023 показва обратното за Port Scan (Recon-PortScan): таванът там е около 85\%. Така качеството на една IDS се определя преди всичко от това дали признаците на потока носят отпечатъка на атаката.

\medskip
\noindent\textsf{\textbf{\color{navy}Ключови думи:}} откриване на прониквания, център за данни, облак, CSE-CIC-IDS2018, DDoS, Brute Force, Botnet, Random Forest, SVM, XGBoost, подбор на признаци, преобучение, bootstrap, тест на McNemar.
